\label{sec:conclusion}

Algorithmic redistricting is implementable through three distinct pathways:
commission adoption, statutory enactment, and court-ordered remedy.  Each
pathway places different institutional actors in the role of map adopter, but
all three share a common structure: the algorithm produces a starting plan
that satisfies quantifiable mathematical criteria, and human oversight bodies
apply judgment to the residual questions (communities of interest, VRA
specifics) that the algorithm cannot resolve without partisan data.

The Arizona case study illustrates that the commission pathway is legally
established and operationally feasible.  The AIRC's constitutional authority
is confirmed; its five-member bipartisan structure provides built-in balance;
and its existing community-of-interest hearing process can absorb the algorithm
output as a constrained starting plan without fundamental change to the
commission's role.

The California case study illustrates that the statute pathway is feasible
where a state's redistricting criteria map cleanly onto the algorithm's weight
parameters.  California's criteria---population equality, VRA compliance, county
contiguity, compactness, communities of interest---correspond to documented
features of the \texttt{redist} system.  The primary implementation challenge
is calibrating the county-sticky weight to California's specific county
structure, which varies enormously in size from Los Angeles County to
Alpine County.

The North Carolina case study illustrates that the court-order pathway is the
most legally stable of the three, because it places the burden of demonstrating
constitutional defect on challengers rather than defenders.  The special master
structure is well-established in federal redistricting practice and available
to state courts.  The algorithm's inability to access partisan data eliminates
the most common basis for constitutional challenge---partisan intent---leaving
challengers with the more demanding task of demonstrating constitutional
violation through effects alone.

Common to all three pathways is the need for data preparation, public
transparency, and a structured challenge process.  These challenges are
real but tractable.  The \texttt{redist} system's documented inputs, reproducible
outputs, and quantitative criteria-compliance reports address transparency and
challenge-process concerns directly.  Data preparation, while time-consuming
for large states, is a one-time cost that supports subsequent redistricting
cycles.

Adoption of algorithmic redistricting at scale would likely follow the pattern
of any significant institutional reform: early adopters in states with
independent commissions or active state-court supervision, followed by
broader uptake as the early cases establish legal precedent and operational
best practices.  The three pathways analyzed here provide the institutional
templates for that process.
